\documentclass[10pt,twocolumn]{article}
\usepackage[T1]{fontenc}
\usepackage[utf8]{inputenc}
\usepackage{lmodern}
\usepackage[margin=1.8cm,top=2.4cm,bottom=2cm,columnsep=0.6cm]{geometry}
\usepackage{fancyhdr}
\usepackage{titlesec}
\usepackage{booktabs}
\usepackage{amsmath,amssymb,amsfonts}
\usepackage{microtype}
\usepackage{xcolor}
\usepackage{mdframed}
\usepackage{parskip}
\usepackage{enumitem}
\usepackage{hyperref}

\definecolor{rulered}{RGB}{180,30,30}
\definecolor{boxgray}{RGB}{245,245,245}
\definecolor{keywordgray}{RGB}{100,100,100}

\pagestyle{fancy}
\fancyhf{}
\fancyhead[L]{\textsc{\small Claude Mag}}
\fancyhead[C]{\small\textit{Machine Learning Research Digest}}
\fancyhead[R]{\small 2026-08-19}
\fancyfoot[C]{\small\thepage}
\renewcommand{\headrulewidth}{0.8pt}

\titleformat{\section}{\normalsize\bfseries\color{rulered}}{}{0pt}{}%
  [\vspace{-4pt}\color{rulered}\rule{\columnwidth}{0.4pt}]
\titlespacing{\section}{0pt}{8pt}{4pt}

\hypersetup{colorlinks=true,urlcolor=rulered,linkcolor=black}

\begin{document}
\setlength{\parindent}{0pt}
\setlength{\parskip}{4pt}

{\small\color{keywordgray}\textbf{Keywords:}
  flow matching \textbullet{}
  efficient generative models \textbullet{}
  image synthesis \textbullet{}
  multidimensional conditioning}\par\vspace{4pt}

{\LARGE\bfseries\raggedright
  XYZFlow: Scaling Multidimensional Shortcut
  Flows for Efficient Generative Modeling}\par\vspace{4pt}

{\large\itshape\raggedright
  Adding Temporal History and Spatial Context to Flow Matching
  Straightens Trajectories, Enabling High-Quality Generation
  in Fewer Steps Without a Teacher Model}\par\vspace{6pt}

{\small\textbf{By} Jinxiu Liu, Xuanming Liu, Kangfu Mei,
  Yandong Wen, Weiyang Liu}\par\vspace{2pt}

{\small\color{keywordgray}arXiv:2608.12276 \textbullet{} Preprint, August 2026 \textbullet{}
  CUHK; Westlake University; Johns Hopkins University}
\par\vspace{2pt}\noindent{\color{rulered}\rule{\columnwidth}{0.6pt}}\vspace{6pt}

Diffusion models produce state-of-the-art images but require dozens of iterative
sampling steps. The standard acceleration path --- distilling a slow teacher into
a few-step student --- inherits the teacher's quality ceiling and demands careful
training. XYZFlow takes a different route: it redesigns flow matching itself by
conditioning probability paths on denoising history (temporal dimension Y) and
neighboring spatial patches (spatial dimension Z), in addition to the standard
denoising time (dimension X). Richer conditioning straightens the flow trajectory,
enabling high-quality generation in far fewer steps with no teacher model required.

\begin{mdframed}[backgroundcolor=boxgray,linecolor=rulered,linewidth=1pt,
    innerleftmargin=6pt,innerrightmargin=6pt,
    innertopmargin=6pt,innerbottommargin=6pt]
\textbf{\small AT A GLANCE}\par\vspace{4pt}
\begin{description}[leftmargin=0pt,labelsep=4pt,itemsep=2pt,parsep=0pt,
    font=\small\bfseries]
  \item[Problem:] Flow/diffusion models require many sampling steps for high quality;
    distillation-based speedups depend on teacher quality and are hard to train.
  \item[Why it matters:] Efficient generation is critical for interactive and
    high-throughput image synthesis applications.
  \item[Method:] Extend flow matching to three conditioning dimensions: denoising
    time (X), denoising history (Y), and spatial patch context (Z).
  \item[Result:] Higher quality per function evaluation than distilled few-step
    models; no teacher model required; quality improves monotonically with conditioning.
  \item[Evidence:] FID and CLIP-score improvements over 1--4 step baselines across
    multiple image benchmarks.
\end{description}
\end{mdframed}

\section{The Big Picture}

Flow matching learns a continuous-time transformation that maps noise
$z \sim \mathcal{N}(0, I)$ to data $x \sim p_{\text{data}}$ via a
probability path $(p_t)_{t \in [0,1]}$. Standard flow matching parameterizes
this path along a single temporal dimension $t$ (the denoising time), yielding
a velocity field $v_\theta(z, t)$ that must integrate to transform noise into
data in a finite number of steps.

The core challenge is trajectory curvature: curved probability paths require
many small integration steps for accurate traversal. Straight paths can be
traversed in fewer steps with larger step sizes. Existing efficient methods
reduce curvature after the fact via distillation, but this is computationally
expensive and inherits the teacher model's limitations.

XYZFlow introduces two additional conditioning dimensions that straighten the
path \emph{by design} during training: the denoising history (accumulated context
from prior steps) and the spatial context (generated patches in neighboring regions).
By conditioning on more of the world state at each step, the model's uncertainty
about the optimal next step decreases, and the induced trajectory becomes straighter.

\section{Why Existing Approaches Fall Short}

\textbf{One/two-step distillation} (Consistency Models, LCM, TurboFlow) trains a
student model to match the teacher's output in one or two steps. Quality is bounded
by the teacher, and the distillation objective is sensitive to hyperparameters. The
student also cannot generalize beyond the teacher's mode coverage.

\textbf{Adversarial diffusion distillation} (ADD, DMD) adds a discriminator to
the distillation loss, improving one-step quality but further complicating training
and introducing GAN instabilities.

\textbf{Autoregressive generation} achieves efficiency through sequential patch
prediction but sacrifices the parallelism of diffusion-based approaches and requires
discrete tokenization.

XYZFlow achieves the efficiency gains of distilled models without a teacher,
and the expressivity gains of autoregressive models without discrete tokenization.

\section{Core Mechanism}

XYZFlow's velocity field is conditioned on three quantities at each step:

\begin{itemize}[itemsep=2pt,parsep=0pt]
  \item \textbf{X --- denoising time $t$:} Standard conditioning on current
    noise level, shared with all flow matching models.
  \item \textbf{Y --- denoising history $h_t = (z_{t_1}, z_{t_2}, \ldots)$:}
    A memory of the partial trajectory, conditioning on states at a set of
    prior denoising times $t_1 > t_2 > \cdots > t$. This non-Markovian
    conditioning allows the model to exploit accumulated context about
    the direction of the trajectory.
  \item \textbf{Z --- spatial context $s_i$:} During spatial dimension
    generation (sequential patch prediction), conditioning on the set
    of already-generated patches $s_{<i}$ provides spatial consistency
    that reduces uncertainty about each new patch's target state.
\end{itemize}

The velocity field is:
\[
  v_\theta(z_t, t, h_t, s_{<i}) \approx \nabla_{z_t} \log p_t(z_t \mid h_t, s_{<i})
\]
with richer conditioning providing a more informative gradient direction at each step.

\section{Technical Formulation}

Standard flow matching minimizes the conditional flow matching loss:
\[
  \mathcal{L}_{\text{FM}}(\theta) = \mathbb{E}_{t, z_0, z_1}
  \left[\|v_\theta(z_t, t) - (z_1 - z_0)\|^2\right]
\]
where $z_t = (1-t)z_0 + tz_1$ is the linear interpolation path.

XYZFlow extends this to the multidimensional conditioned loss:
\[
  \mathcal{L}_{\text{XYZ}}(\theta) = \mathbb{E}_{t, z_0, z_1, h_t, s_{<i}}
  \left[\|v_\theta(z_t, t, h_t, s_{<i}) - (z_1 - z_0)\|^2\right]
\]

The key theoretical result is that the \emph{optimal} velocity $v^*(z_t, t, h_t, s_{<i})$
is straighter (lower curvature) than $v^*(z_t, t)$ whenever the conditioning
variables $h_t$ and $s_{<i}$ are informative about the target $z_1$.
This follows from the variance decomposition:
\[
  \mathbb{E}[\|z_1 - z_0\|^2] \geq \mathbb{E}[\|z_1 - \mathbb{E}[z_1 \mid h_t, s_{<i}, z_t]\|^2]
\]
The residual uncertainty under conditioning is lower, implying the conditioned
target is closer to the true target at each step.

\section{Learning or Inference Procedure}

Training follows standard flow matching: sample pairs $(z_0, z_1)$, sample
time $t$, compute $z_t$, retrieve history $h_t$ from prior steps of the
current trajectory, retrieve spatial context $s_{<i}$ from previously
generated patches, and minimize $\mathcal{L}_{\text{XYZ}}$.

At inference, generation proceeds in two nested loops:
\begin{enumerate}[itemsep=2pt,parsep=0pt]
  \item \textbf{Temporal loop:} Take $N_T$ denoising steps ($N_T = 2$--$8$
    for efficient generation, vs.\ $50$--$1000$ for standard diffusion).
    Maintain history $h_t$ across steps.
  \item \textbf{Spatial loop:} At each denoising step, generate patches
    sequentially, conditioning each patch on the already-generated patches.
    $N_S$ patches are generated per step.
\end{enumerate}

\section{What the Guarantee Says}

The paper proves that, under perfect conditioning, the XYZFlow velocity field
induces a straighter trajectory than standard flow matching, requiring fewer
integration steps for the same approximation error. The practical guarantee is:
for a fixed number of function evaluations (NFE), XYZFlow achieves lower FID
than unconditioned flow matching and distilled models at matched NFE.

\section{Experimental Findings}

\begin{itemize}[itemsep=2pt,parsep=0pt]
  \item \textbf{Quality vs.\ NFE:} At NFE$=2$, XYZFlow outperforms LCM and
    TurboFlow by 3--5 FID points on ImageNet-256 and COCO.
  \item \textbf{No teacher required:} XYZFlow trained from scratch matches
    or exceeds the quality of distilled models trained from a strong teacher.
  \item \textbf{Monotonic improvement:} Adding Y conditioning over X improves
    FID by 2--3 points; adding Z conditioning further improves by 1--2 points.
    X $<$ XY $<$ XYZ holds consistently across NFE settings.
  \item \textbf{Scaling:} Quality scales gracefully with NFE; there is no
    saturation plateau seen in distilled models at 2--4 steps.
\end{itemize}

\section{Ablations and Interpretation}

\textbf{History length:} Conditioning on 2--3 prior states (vs.\ just $t-1$)
captures most of the benefit; conditioning on the full history provides
marginal additional gain at higher memory cost.

\textbf{Spatial patch order:} Raster-scan order outperforms random order,
confirming that spatial locality in the conditioning signal matters.

\textbf{Architecture:} Cross-attention over history states outperforms
concatenation for integrating the Y and Z conditioning signals.

\section*{Reference}

Jinxiu Liu, Xuanming Liu, Kangfu Mei, Yandong Wen, Weiyang Liu.
\textbf{XYZFlow: Scaling Multidimensional Shortcut Flows for Efficient Generative Modeling}.
arXiv:2608.12276, August 2026.
\url{https://arxiv.org/abs/2608.12276}

\end{document}
